\documentclass{article}

% Language setting
% Replace `english' with e.g. `spanish' to change the document language
\usepackage[english]{babel}

% Set page size and margins
% Replace `letterpaper' with `a4paper' for UK/EU standard size
\usepackage[letterpaper,top=2cm,bottom=2cm,left=3cm,right=3cm,marginparwidth=1.75cm]{geometry}

% Useful packages
\usepackage{amsmath}
\usepackage{graphicx}
\usepackage[colorlinks=true, allcolors=blue]{hyperref}

\title{A Reduced Multicompartment Model for Pulsation- and AQP4-Mediated Water Exchange}
\author{You}

\begin{document}
\maketitle

\begin{abstract}
Your abstract.
\end{abstract}

 

We consider a reduced model of the periarteriolar glymphatic interface consisting of four compartments:
the perivascular space (PVS), denoted by \(P\), the extracellular space/interstitial fluid compartment,
denoted by \(E\), the astrocytic endfoot compartment, denoted by \(A\), and the astrocytic process/soma
compartment, denoted by \(S\). The model is intended to describe how arterial radius oscillations,
inter-endfoot gap dynamics, and aquaporin-4 (AQP4)-mediated water exchange regulate fluid exchange
between the PVS and the ECS.

The vessel radius is prescribed as an external input \(R_v(t)\), representing cardiac pulsation,
functional hyperemia, or sleep-like vasomotion. This choice follows previous poroelastic models in
which vessel wall motion is prescribed and the resulting PVS pressure and fluid exchange are computed.
Kedarasetti et al. showed that temporally asymmetric arteriolar dilation can drive net PVS-to-ECS
fluid penetration, whereas Causemann et al. emphasized that, at the local gliovascular interface,
pressure-driven exchange occurs primarily through inter-endfoot gaps rather than through the
AQP4-rich endfoot membrane. 


\section{Vessel-radius forcing and PVS compression}

Let \(R_v(t)\) be the radius of the penetrating arteriole and let \(R_o(t)\) be the inner radius of the
astrocytic endfoot sheath. The geometric volume of the PVS over a vessel segment of length \(L\) is
approximated by
\begin{equation}
    V_P^{\mathrm{geom}}(t)
    =
    \pi L \left( R_o(t)^2 - R_v(t)^2 \right).
\end{equation}
To account for the mechanical coupling between the expanding vessel and the endfoot sheath, we set
\begin{equation}
    R_o(t)
    =
    R_{o0}
    +
    \eta \left( R_v(t)-R_{v0} \right),
\end{equation}
where \(0\leq \eta\) measures the degree to which the endfoot sheath follows the vascular motion.
When \(\eta=0\), the outer PVS boundary is fixed and vessel dilation compresses the PVS. When
\(\eta\) is large, tangential stretch of the endfoot sheath can partially compensate or even overcome
radial PVS compression.

The geometric source term generated by vessel motion is defined as
\begin{equation}
    S_R(t)
    =
    -\frac{dV_P^{\mathrm{geom}}}{dt}.
\end{equation}
Thus, when vessel dilation decreases the available PVS volume,
\[
    \frac{dV_P^{\mathrm{geom}}}{dt}<0,
    \qquad
    S_R(t)>0,
\]
which acts as a pressure source in the PVS.


\section{Prescribed vessel-radius waveforms}

We consider three classes of vessel-radius waveforms.

\paragraph{Cardiac pulsation.}
A small-amplitude, high-frequency oscillation is represented by
\begin{equation}
    R_v(t)
    =
    R_{v0}
    \left[
    1+\epsilon_c \sin(2\pi f_c t)
    \right],
\end{equation}
where \(f_c\) is the cardiac frequency and \(\epsilon_c\) is the relative amplitude.

\paragraph{Symmetric dilation.}
A symmetric dilation event is represented by
\begin{equation}
    R_v(t)
    =
    R_{v0}
    \left[
    1+
    A\exp\left(-\frac{(t-t_0)^2}{2\sigma^2}\right)
    \right].
\end{equation}

\paragraph{Asymmetric functional hyperemia.}
A fast dilation followed by a slower return to baseline is represented by
\begin{equation}
    R_v(t)
    =
    R_{v0}
    \left[
    1+A f(t)
    \right],
\end{equation}
where
\begin{equation}
    f(t)
    =
    \begin{cases}
    1-\exp(-t/\tau_{\mathrm{up}}), & 0<t<t_p, \\[4pt]
    \exp(-(t-t_p)/\tau_{\mathrm{down}}), & t\geq t_p,
    \end{cases}
\end{equation}
with
\[
    \tau_{\mathrm{up}}<\tau_{\mathrm{down}}.
\]
This waveform captures the characteristic rapid onset and slow recovery of functional hyperemia.


\section{Water fluxes}

The main PVS-to-ECS exchange pathway is assumed to be the inter-endfoot gap. We model the gap flow by
\begin{equation}
    Q_{PE}
    =
    G_{\mathrm{gap}}(w)(p_P-p_E),
\end{equation}
where \(p_P\) and \(p_E\) are the hydrostatic pressures in the PVS and ECS, respectively. The gap
conductance is assumed to depend sensitively on the gap width \(w(t)\):
\begin{equation}
    G_{\mathrm{gap}}(w)
    =
    G_{\mathrm{gap},0}
    \left(
    \frac{w}{w_0}
    \right)^m,
    \qquad m=2 \text{ or }3.
\end{equation}

AQP4-mediated water exchange across the PVS-facing endfoot membrane is modeled by a Starling-type law.
Taking positive flux from the PVS into the endfoot, we define
\begin{equation}
    Q_{PA}
    =
    L_{\mathrm{AQP4}}
    \left[
    (p_P-p_A)
    -
    RT(\Pi_P-\Pi_A)
    \right],
\end{equation}
where \(p_A\) is the endfoot pressure, \(\Pi_P\) and \(\Pi_A\) are the effective osmolarities in the PVS
and endfoot, and \(L_{\mathrm{AQP4}}\) is the AQP4-controlled hydraulic permeability.

The water exchange between the endfoot and ECS is modeled as
\begin{equation}
    Q_{AE}
    =
    L_{AE}
    \left[
    (p_A-p_E)
    -
    RT(\Pi_A-\Pi_E)
    \right],
\end{equation}
where \(L_{AE}\) may be smaller than \(L_{\mathrm{AQP4}}\), reflecting a structural asymmetry between the
PVS-facing and ECS-facing sides of the endfoot.

Finally, water redistribution between the endfoot and the astrocytic process/soma compartment is given by
\begin{equation}
    Q_{AS}
    =
    G_{AS}
    \left[
    (p_A-p_S)
    -
    RT(\Pi_A-\Pi_S)
    \right].
\end{equation}

\section{Endfoot-gap dynamics}

We hypothesize that endfoot deformation dynamically modulates the inter-endfoot gap width. The preferred
gap width is modeled as
\begin{equation}
    w_*(t)
    =
    w_0
    +
    \chi_R\left(R_v(t)-R_{v0}\right)
    +
    \chi_A\left(V_A(t)-V_{A0}\right).
\end{equation}
Here \(\chi_R>0\) represents tangential stretching of the endfoot sheath induced by vessel dilation,
while \(\chi_A\) determines how endfoot volume changes affect the gap. If \(\chi_A>0\), endfoot swelling
is associated with tangential footprint shrinkage and gap opening. If \(\chi_A<0\), endfoot swelling
compresses the ECS/gap and decreases gap conductance.

To allow for a phase lag between vessel motion, endfoot volume change, and gap opening, we use the
relaxation equation
\begin{equation}
    \tau_w \frac{dw}{dt}
    =
    -\left(w-w_*(t)\right),
\end{equation}
where \(\tau_w\) is the characteristic time scale of gap adjustment.
\section{Compartment pressure and volume equations}

The PVS pressure evolves according to
\begin{equation}
    C_P \frac{dp_P}{dt}
    =
    S_R(t)
    -
    Q_{PE}
    -
    Q_{PA}
    -
    Q_{P,out}
    +
    Q_{P,in},
\end{equation}
where \(C_P\) is the PVS compliance. The downstream PVS outflow is modeled by
\begin{equation}
    Q_{P,out}
    =
    G_{P,out}(p_P-p_{out}).
\end{equation}

The ECS pressure evolves according to
\begin{equation}
    C_E \frac{dp_E}{dt}
    =
    Q_{PE}
    +
    Q_{AE}
    -
    Q_{E,out},
\end{equation}
where
\begin{equation}
    Q_{E,out}
    =
    G_{E,out}(p_E-p_{out,E})
\end{equation}
represents drainage from the ECS into downstream clearance pathways.

The endfoot volume satisfies
\begin{equation}
    \frac{dV_A}{dt}
    =
    Q_{PA}
    -
    Q_{AE}
    -
    Q_{AS}.
\end{equation}
The endfoot pressure is related to its volume by a compliance law:
\begin{equation}
    p_A
    =
    p_{A0}
    +
    \frac{V_A-V_{A0}}{C_A}.
\end{equation}

The astrocytic process/soma compartment is modeled as
\begin{equation}
    \frac{dV_S}{dt}
    =
    Q_{AS}
    -
    Q_{S,out}.
\end{equation}
In the simplest version, this compartment can be treated as a large reservoir with
\[
    p_S=p_{S0},
    \qquad
    \Pi_S=\Pi_{S0}.
\]
\subsection{Compartment pressure and volume equations}

The PVS pressure evolves according to
\begin{equation}
    C_P \frac{dp_P}{dt}
    =
    S_R(t)
    -
    Q_{PE}
    -
    Q_{PA}
    -
    Q_{P,out}
    +
    Q_{P,in},
\end{equation}
where \(C_P\) is the PVS compliance. The downstream PVS outflow is modeled by
\begin{equation}
    Q_{P,out}
    =
    G_{P,out}(p_P-p_{out}).
\end{equation}

The ECS pressure evolves according to
\begin{equation}
    C_E \frac{dp_E}{dt}
    =
    Q_{PE}
    +
    Q_{AE}
    -
    Q_{E,out},
\end{equation}
where
\begin{equation}
    Q_{E,out}
    =
    G_{E,out}(p_E-p_{out,E})
\end{equation}
represents drainage from the ECS into downstream clearance pathways.

The endfoot volume satisfies
\begin{equation}
    \frac{dV_A}{dt}
    =
    Q_{PA}
    -
    Q_{AE}
    -
    Q_{AS}.
\end{equation}
The endfoot pressure is related to its volume by a compliance law:
\begin{equation}
    p_A
    =
    p_{A0}
    +
    \frac{V_A-V_{A0}}{C_A}.
\end{equation}

The astrocytic process/soma compartment is modeled as
\begin{equation}
    \frac{dV_S}{dt}
    =
    Q_{AS}
    -
    Q_{S,out}.
\end{equation}
In the simplest version, this compartment can be treated as a large reservoir with
\[
    p_S=p_{S0},
    \qquad
    \Pi_S=\Pi_{S0}.
\]

\section{Tracer transport and clearance}

Let \(c_P(t)\) and \(c_E(t)\) denote the tracer or waste concentrations in the PVS and ECS. We define
\[
    Q_{PE}^{+}=\max(Q_{PE},0),
    \qquad
    Q_{EP}^{+}=\max(-Q_{PE},0).
\]
The tracer balance in the PVS is
\begin{equation}
    \frac{d}{dt}(V_P c_P)
    =
    Q_{P,in}c_{in}
    -
    Q_{PE}^{+}c_P
    +
    Q_{EP}^{+}c_E
    -
    Q_{P,out}c_P.
\end{equation}
The tracer balance in the ECS is
\begin{equation}
    \frac{d}{dt}(V_E c_E)
    =
    Q_{PE}^{+}c_P
    -
    Q_{EP}^{+}c_E
    -
    Q_{E,out}c_E
    -
    k_{\mathrm{deg}}V_Ec_E
    +
    S_c.
\end{equation}

The clearance performance can be quantified by the decay half-time of \(c_E(t)\), the cumulative
downstream removal
\begin{equation}
    \mathcal{C}_{out}(T)
    =
    \int_0^T Q_{E,out}(t)c_E(t)\,dt,
\end{equation}
and the cycle-averaged PVS-to-ECS exchange
\begin{equation}
    \overline{Q}_{PE}
    =
    \frac{1}{T}
    \int_0^T Q_{PE}(t)\,dt.
\end{equation}
\end{document}
